%% Subresultants in roots without denominators
%%
%% Generated by Claude Opus 5 (Anthropic) between 15 and 17 September 2026.
%% Published unverified in AI-Generated Mathematical Explorations, exploration 003.
%%
%% This file is the AI output as generated, with four editorial changes, all of them
%% made by the curator and listed here in full:
%%   1. this header;
%%   2. a provenance note added after \maketitle (the \fbox block below), so that
%%      the PDF carries its own status when it travels alone; the line
%%      "% TODO: autores y afiliacion" of the original was removed with it, and
%%      urlcolor was added to the hyperref options so that the link inside it is
%%      coloured like the other links;
%%   3. the date, left as \today in the original, was fixed to the date on
%%      which the note was generated, so that recompiling does not change it;
%%   4. the framing of the abstract, of the paragraph of the introduction that
%%      states the questions and of the first item of the list of results was
%%      rewritten, to say plainly that the ratio is a polynomial in the roots for
%%      trivial reasons -- it is the subresultant itself -- and that the content
%%      of the work is the explicit, quotient-free form of it; the word "paper"
%%      was replaced by "note" throughout.  No statement, proof, formula or
%%      example was touched.
%% Apart from those, the text is as generated, including the "% TODO" comment
%% before the bibliography, which is the AI's own note that it wrote the
%% bibliographic data from memory.  The original file is reproduced verbatim in
%% ai-output.md, next to this one.

\documentclass[11pt,reqno]{amsart}

\usepackage[T1]{fontenc}
\usepackage[utf8]{inputenc}
\usepackage{amsmath,amssymb,amsthm,mathtools}
\usepackage{array}
\usepackage[margin=1.15in]{geometry}
\usepackage[colorlinks=true,linkcolor=blue,citecolor=blue,urlcolor=blue]{hyperref}

\theoremstyle{plain}
\newtheorem{theorem}{Theorem}[section]
\newtheorem{proposition}[theorem]{Proposition}
\newtheorem{lemma}[theorem]{Lemma}
\newtheorem{corollary}[theorem]{Corollary}
\theoremstyle{definition}
\newtheorem{definition}[theorem]{Definition}
\newtheorem{remark}[theorem]{Remark}
\newtheorem{example}[theorem]{Example}

\DeclareMathOperator{\Res}{Res}
\DeclareMathOperator{\sres}{sres}
\DeclareMathOperator{\sgn}{sgn}
\DeclareMathOperator{\sg}{sg}
\newcommand{\bA}{\mathbb{A}}
\newcommand{\bX}{\Xi}
\newcommand{\cO}{\mathcal{O}}
\newcommand{\cS}{\mathcal{S}}
\newcommand{\cV}{\mathcal{V}}
\newcommand{\cC}{\mathcal{C}}
\newcommand{\ZZ}{\mathbb{Z}}
\newcommand{\QQ}{\mathbb{Q}}
\newcommand{\rev}{\mathrm{rev}}

\title[Subresultants in roots without denominators]
      {Subresultants in roots without denominators}
\author{}
\date{September 15, 2026}

\begin{document}

\begin{abstract}
The formulas of D'Andrea, Krick and Szanto express subresultants of univariate
polynomials as a ratio of two determinants, the denominator being a Vandermonde
determinant in the roots. That this ratio is a polynomial is not the issue: it is
the subresultant itself, up to a power of the leading coefficient of $g$, and the
division is exact already over $\ZZ$. What is missing is a workable expression for
it, and that is what we give.
The quotient is a multi-Schur function in the sense of Lascoux, built from two
alphabets: the roots $\bX$ of one polynomial and the difference $\bX-\bA$ of the
two root alphabets. We give three explicit forms for it: a Schur expansion whose
coefficients are maximal minors of the band matrix of $f$; a Jacobi--Trudi
determinant with one alphabet per row, with the sign determined exactly; and an
expansion in the basis $\{s_\nu(\bX-\bA)\}$ whose coefficients are \emph{skew}
Schur functions of the roots of $f$, with an explicit vanishing criterion.
For the principal coefficient of the classical subresultant the quotient
collapses to a single supersymmetric Schur function indexed by the rectangle
$(d_1-k)\times(d_2-k)$, recovering a theorem of Lascoux and Pragacz, and
equivalently to a Toeplitz determinant in the Taylor coefficients of $f/g$,
which is the Pad\'e approximation form. We also show that the quotient is the
determinant of divided differences of the rows, which settles the case of
multiple roots at once. Finally we show that the naive multivariate analogue of
the divisibility statement is false, and we identify the structural reason.
\end{abstract}

\maketitle

\begin{center}\small
\fbox{\begin{minipage}{0.92\textwidth}
\textbf{Provenance.} This note was generated by an AI system (Claude Opus 5,
Anthropic) between 15 and 17 September 2026, in the context of research
activities of Santiago Laplagne at the Instituto de C\'alculo, Universidad de
Buenos Aires. It is published \textbf{unverified}: no human author claims these
mathematical results, and the arguments have not been independently checked.
Current status, original AI output, and how to report an error or a
verification: \url{https://github.com/slap/ai-mathematical-explorations}
\end{minipage}}
\end{center}

\section{Introduction}

Let $f,g$ be univariate polynomials of degrees $d_1,d_2$ over a field $K$, and
let $\xi_1,\dots,\xi_{d_2}$ be the roots of $g$ in $\overline{K}$. The classical
Poisson product formula
\[
  \Res(f,g)=(-1)^{d_1d_2}b_{d_2}^{d_1}\prod_{j=1}^{d_2}f(\xi_j)
\]
expresses the resultant as a polynomial expression in the roots of $g$ and the
coefficients of $f$, and a usable one. For subresultants the known analogues are
not of that shape:
Hong's formula \cite{Hong}, and more generally the formulas of D'Andrea, Krick
and Szanto \cite{DKS05}, express subresultants as a \emph{ratio} of two
determinants, the denominator being the Vandermonde determinant
$\cV_{d_2}=\prod_{i<j}(\xi_j-\xi_i)$ of the roots of $g$. Concretely,
\cite[Thm.~2.2]{DKS05} states
\begin{equation}\label{eq:dks}
  \Delta_\cS \;=\; \sg(\cS)\, b_{d_2}^{\,t^*-d_2+1}\,
  \frac{|\cO_\cS|}{\cV_{d_2}},
\end{equation}
where $\cO_\cS$ is a $d_2\times d_2$ matrix whose columns are indexed by the
roots of $g$ (all notation is recalled in Section~\ref{sec:setup}).

That the right-hand side of \eqref{eq:dks} is a polynomial in the roots of $g$ and
the coefficients of $f$ needs no proof: it is the subresultant, up to the factor
$\sg(\cS)b_{d_2}^{\,t^*-d_2+1}$, and subresultants are polynomials in the
coefficients of $f$ and $g$. The formula, however, does not exhibit it as one, and
a quotient is not something one can compute with, expand, specialize or bound.
This note therefore asks for the polynomial itself: what is
$|\cO_\cS|/\cV_{d_2}$ explicitly, in a form free of denominators, and in which
basis of symmetric functions is it naturally expressed? Is the division exact over
$\ZZ$? Do the answers extend to multiple roots and to the multivariate setting?

\medskip
\noindent\textbf{Results.} Write
\[
  P_\cS \;:=\; \frac{|\cO_\cS|}{\cV_{d_2}} .
\]
\begin{itemize}
\item \textbf{Divisibility over $\ZZ$} (Theorem~\ref{thm:div}). This is the
elementary step, and it is where the denominator goes. Each column of $\cO_\cS$
depends on a single root, so $|\cO_\cS|$ vanishes on every diagonal
$\xi_i=\xi_j$; since the roots of a generic $g$ are algebraically independent,
$\cV_{d_2}$ divides $|\cO_\cS|$ in $\ZZ[a_0,\dots,a_{d_1}][\xi_1,\dots,\xi_{d_2}]$
and $P_\cS$ is symmetric in the roots. The whole content of the problem is
therefore the identification of $P_\cS$.

\item \textbf{Schur expansion} (Theorem~\ref{thm:cb}). The matrix $\cO_\cS$
factors as $\rho\cdot W$, where $W$ is the evaluation matrix of monomials at the
roots and $\rho$ is the coefficient matrix of the rows. The Cauchy--Binet formula
and the bialternant formula give
$P_\cS=\sum_E \det(\rho_{\bullet,E})\,s_{\lambda(E)}(\bX)$: the quotient expands
in the Schur basis of the roots of $g$ with coefficients the maximal minors of the
band matrix of $f$.

\item \textbf{Multi-Schur form, with the exact sign} (Theorem~\ref{thm:multischur}).
$P_\cS$ is a Jacobi--Trudi determinant with \emph{one alphabet per row}: the
alphabet $\bX$ for the monomial rows and the difference alphabet $\bX-\bA$ for
the rows coming from $f$. The global sign is $(-1)^{\binom{d_2}{2}}\sgn(\sigma)$,
and the statement holds for any ordering of the rows.

\item \textbf{The classical case is a rectangle} (Theorem~\ref{thm:rect}). When
$\cS=\{1,x,\dots,x^{k-1}\}$, that is for the principal coefficient of the
classical subresultant, the determinant collapses and
$P_{\cS_k}=a_{d_1}^{d_2-k}\,s_{((d_1-k)^{d_2-k})}(\bX-\bA)$: a single
supersymmetric Schur function indexed by the rectangle
$(d_1-k)\times(d_2-k)$. This recovers a theorem of Lascoux and Pragacz
\cite{LP} from \eqref{eq:dks}. The collapse uses in an essential way that the
exponents of $\cS$ are \emph{consecutive}.

\item \textbf{Toeplitz / Pad\'e form} (Theorem~\ref{thm:toeplitz}). Equivalently,
the principal coefficient is a Toeplitz determinant in the Taylor coefficients of
$f^\rev/g^\rev$, which is the expansion of $f/g$ at infinity. This is the
subresultant--Pad\'e dictionary.

\item \textbf{The natural basis for arbitrary $\cS$} (Theorem~\ref{thm:skew}).
For general $\cS$ the quotient is \emph{not} a single supersymmetric Schur
function; the obstruction is caused by gaps in the exponent set. It does expand in
that basis, by the Cauchy--Binet argument applied in the difference alphabet, and
the resulting expansion is dual to the one of Theorem~\ref{thm:cb}: the
coefficients are \emph{skew} Schur functions $s_{\lambda/\mu}(\bA)$ of the roots
of $f$ alone, and they are nonzero exactly when $\mu\subseteq\lambda$ and
$\lambda/\mu$ has no column longer than $d_1$.

\item \textbf{Multiple roots} (Theorem~\ref{thm:newton}). The quotient equals the
determinant of the \emph{divided differences} of the row polynomials,
$P_\cS=\det\big(R_i[\xi_1,\dots,\xi_j]\big)$. This form is manifestly polynomial,
and specializing repeated roots turns divided differences into derivatives; so the
confluent case requires no new formula.

\item \textbf{The multivariate case} (Proposition~\ref{prop:multi}). The naive
analogue of Theorem~\ref{thm:div} is \emph{false}. Already in the smallest
example of \cite[Ex.~3.4]{DKS05} the generalized Vandermonde determinant $\cV_T$
divides none of the alternants occurring in the numerator. The reason is
structural: the proof of Theorem~\ref{thm:div} uses that alternating polynomials
in one variable per point form a free module of rank one over the symmetric ones,
and this fails for points of $\mathbb{A}^n$ with $n\ge 2$.
\end{itemize}

\medskip
\noindent\textbf{Status of the statements.} Theorems~\ref{thm:div},
\ref{thm:cb}, \ref{thm:multischur}, \ref{thm:rect}, \ref{thm:toeplitz},
\ref{thm:skew} and \ref{thm:newton} are proved. All of them, together with
Proposition~\ref{prop:multi}, have in addition been verified by symbolic
computation over a range of parameters; Section~\ref{sec:verif} reports the
scope of those computations, which were also used to discover the statements.

\section{Setup and notation}\label{sec:setup}

We follow \cite[\S2]{DKS05}. Let
\[
  f=\sum_{l=0}^{d_1}a_lx^l,\qquad g=\sum_{l=0}^{d_2}b_lx^l
\]
be generic of degrees $d_1,d_2$, with coefficients algebraically independent over
$\QQ$, and let $K=\QQ(a_0,\dots,a_{d_1},b_0,\dots,b_{d_2})$. Let
$\bX=\{\xi_1,\dots,\xi_{d_2}\}$ be the roots of $g$ and
$\bA=\{\alpha_1,\dots,\alpha_{d_1}\}$ the roots of $f$ in $\overline K$, and put
\[
  \cV_{d_2}:=\det\big(\xi_j^{\,i-1}\big)_{1\le i,j\le d_2}
  =\prod_{1\le i<j\le d_2}(\xi_j-\xi_i).
\]
Fix $t$ with $0\le t\le d_1+d_2-1$ and set
\[
  t^*:=\max\{d_2-1,\,t\},\qquad
  k:=t+1-\max\{0,t-d_1+1\}-\max\{0,t-d_2+1\},
\]
and let $\cS=\{x^{\gamma_1},\dots,x^{\gamma_k}\}\subset K[x]_t$ with
$0\le\gamma_1<\dots<\gamma_k\le t$. Following \cite{DKS05}, $\Delta_\cS$ denotes
the order $t$ subresultant of $f,g$ with respect to $\cS$, and $\sg(\cS)$ the sign
defined there.

\begin{definition}\label{def:rows}
The matrix $\cO_\cS\in\overline K^{\,d_2\times d_2}$ of \cite[Thm.~2.2]{DKS05} has
rows $\big(R_i(\xi_1),\dots,R_i(\xi_{d_2})\big)$, where the row polynomials $R_i$
are
\[
  x^{\gamma_1},\dots,x^{\gamma_k},\quad
  x^{t+1},\dots,x^{t^*},\quad
  f,\ xf,\ \dots,\ x^{t-d_1}f .
\]
We call the first two groups the \emph{monomial rows} and the last one the
\emph{$f$-rows}, and we write
\[
  \Gamma:=\{\gamma_1,\dots,\gamma_k\}\cup\{t+1,\dots,t^*\},\qquad
  \Phi:=\{d_1,d_1+1,\dots,t\}
\]
for the sets of degrees of the two groups ($\Phi=\emptyset$ if $t<d_1$), and
$\varepsilon_i:=\deg R_i$, so that
$\{\varepsilon_i\}_i=\Gamma\sqcup\Phi$ and $\#\Gamma+\#\Phi=d_2$. Finally let
\[
  \rho\in\ZZ[a_0,\dots,a_{d_1}]^{\,d_2\times(t^*+1)},\qquad
  \rho_{i,c}:=[x^c]R_i ,
\]
be the coefficient matrix of the rows: a unit row $e_\gamma$ for each
$\gamma\in\Gamma$, and the shifted rows $(a_0,\dots,a_{d_1})$ of the band matrix
of $f$ for the $f$-rows.
\end{definition}

Throughout we set
\[
  \boxed{\;P_\cS:=\frac{|\cO_\cS|}{\cV_{d_2}}\;}
\]
so that \eqref{eq:dks} reads $\Delta_\cS=\sg(\cS)b_{d_2}^{\,t^*-d_2+1}P_\cS$.

\subsection*{The two alphabets}
We use $\lambda$-ring notation for symmetric functions, as in \cite{Lascoux,LP}.
Let $h_r(\bX)$ denote the complete homogeneous symmetric function, and let
$h_r(\bX-\bA)$ be defined by the generating series
\[
  \sum_{r\ge0}h_r(\bX-\bA)z^r=\frac{\prod_{i=1}^{d_1}(1-\alpha_iz)}
                                    {\prod_{j=1}^{d_2}(1-\xi_jz)} .
\]
Write $f^\rev(z):=\sum_{r=0}^{d_1}a_{d_1-r}z^r$ and
$g^\rev(z):=\sum_{r=0}^{d_2}b_{d_2-r}z^r$, so that
$\prod_j(1-\xi_jz)=g^\rev(z)/b_{d_2}$ and
$\prod_i(1-\alpha_iz)=f^\rev(z)/a_{d_1}$. It is convenient to use the
denominator-free normalizations
\begin{equation}\label{eq:hH}
  h_r:=[z^r]\frac{1}{\prod_j(1-\xi_jz)}=h_r(\bX),
  \qquad
  H_r:=[z^r]\frac{f^\rev(z)}{\prod_j(1-\xi_jz)}=a_{d_1}h_r(\bX-\bA),
\end{equation}
both of which lie in $\ZZ[a][\xi]$, with $h_r=H_r=0$ for $r<0$. We also put
\begin{equation}\label{eq:c}
  c_r:=[z^r]\frac{f^\rev(z)}{g^\rev(z)},\qquad\text{so that}\qquad H_r=b_{d_2}c_r .
\end{equation}
Note that $H_0=a_{d_1}$, \emph{not} $1$; consequently a Jacobi--Trudi determinant
of size $L$ in the $H_r$ carries a factor $a_{d_1}^{L}$, and the size of the
determinant matters.

We grade $\ZZ[a][\xi]$ by
\begin{equation}\label{eq:grading}
  \deg\xi_j=1,\qquad \deg a_{d_1-r}=r .
\end{equation}
With this grading $h_r$ and $H_r$ are homogeneous of degree $r$, the row $R_i$
contributes degree $\varepsilon_i$, hence $|\cO_\cS|$ is homogeneous of degree
$\sum_i\varepsilon_i$ and $\cV_{d_2}$ of degree $\binom{d_2}{2}$.

\section{Divisibility}\label{sec:div}

\begin{lemma}\label{lem:alt}
Let $R$ be an integral domain and $F\in R[\xi_1,\dots,\xi_d]$ such that
$F|_{\xi_i=\xi_j}=0$ for all $i<j$. Then $\prod_{i<j}(\xi_j-\xi_i)$ divides $F$
in $R[\xi_1,\dots,\xi_d]$.
\end{lemma}

\begin{proof}
Since $\xi_j-\xi_i$ is monic of degree $1$ in $\xi_j$, division with remainder by
it is available over any commutative ring: $F=(\xi_j-\xi_i)Q+\varrho$ with
$\varrho$ not involving $\xi_j$ beyond degree $0$ in the division, and evaluating
at $\xi_j=\xi_i$ gives $\varrho=0$. Thus $\xi_j-\xi_i\mid F$ in $R[\xi]$.
Write $F=(\xi_j-\xi_i)Q$. For any other pair $\{k,l\}\ne\{i,j\}$ we have
$0=F|_{\xi_k=\xi_l}=(\xi_j-\xi_i)|_{\xi_k=\xi_l}\cdot Q|_{\xi_k=\xi_l}$, and
$(\xi_j-\xi_i)|_{\xi_k=\xi_l}\ne0$ is a nonzerodivisor of the domain $R[\xi]$, so
$Q$ again vanishes on all the remaining diagonals. Iterating over the
$\binom d2$ pairs gives the claim.
\end{proof}

\begin{theorem}\label{thm:div}
$\cV_{d_2}$ divides $|\cO_\cS|$ in $\ZZ[a_0,\dots,a_{d_1}][\xi_1,\dots,\xi_{d_2}]$.
The quotient $P_\cS$ is symmetric in $\xi_1,\dots,\xi_{d_2}$, homogeneous of degree
$t-d_1+1$ in the $a_l$, and homogeneous of degree
$\sum_i\varepsilon_i-\binom{d_2}{2}$ for the grading \eqref{eq:grading}.
\end{theorem}

\begin{proof}
Each column of $\cO_\cS$ depends on a single root, so specializing $\xi_i=\xi_j$
produces two equal columns and $|\cO_\cS|$ vanishes. The roots of a generic $g$
are algebraically independent over $\QQ$, so $|\cO_\cS|$ is a polynomial in
$\ZZ[a][\xi]$ vanishing on all diagonals, and Lemma~\ref{lem:alt} applies with
$R=\ZZ[a_0,\dots,a_{d_1}]$. Transposing two roots transposes two columns of
$\cO_\cS$ and two factors of $\cV_{d_2}$ with the same sign, so the quotient is
symmetric. The degrees follow from the grading discussion after
\eqref{eq:grading}, each of the $t-d_1+1$ $f$-rows contributing degree $1$ in
the $a_l$.
\end{proof}

\begin{remark}
By \eqref{eq:dks}, $P_\cS=\pm\,b_{d_2}^{-(t^*-d_2+1)}\Delta_\cS$: the quotient
\emph{is} the subresultant, written in the roots of $g$ and the coefficients of
$f$, up to a power of the leading coefficient of $g$. The remaining sections
identify it.
\end{remark}

\section{The Schur expansion}\label{sec:cb}

For $E=\{e_1<\dots<e_{d_2}\}\subseteq\{0,1,\dots,t^*\}$ define the partition
\begin{equation}\label{eq:lambdaE}
  \lambda(E)_r:=e_{d_2+1-r}-(d_2-r),\qquad r=1,\dots,d_2 .
\end{equation}

\begin{theorem}\label{thm:cb}
With the notation of Definition~\ref{def:rows},
\[
  P_\cS=\sum_{E}\det\big(\rho_{\bullet,E}\big)\, s_{\lambda(E)}(\xi_1,\dots,\xi_{d_2}),
\]
the sum being over all $d_2$-subsets $E\subseteq\{0,\dots,t^*\}$. Moreover
$\det(\rho_{\bullet,E})=0$ unless $\Gamma\subseteq E$, in which case
$\det(\rho_{\bullet,E})$ is, up to sign, the maximal minor of the band matrix of
$f$ on the columns $E\smallsetminus\Gamma$.
\end{theorem}

\begin{proof}
Let $W\in\overline K^{(t^*+1)\times d_2}$ be given by $W_{c,j}=\xi_j^{\,c}$. By
Definition~\ref{def:rows}, $\cO_\cS=\rho\cdot W$, and Cauchy--Binet gives
$|\cO_\cS|=\sum_E\det(\rho_{\bullet,E})\det(W_{E,\bullet})$ over $d_2$-subsets
$E$ taken in increasing order. The bialternant formula, in the increasing
convention matching $\cV_{d_2}=\det(\xi_j^{i-1})$, reads
$\det(W_{E,\bullet})=\cV_{d_2}\,s_{\lambda(E)}(\bX)$ with $\lambda(E)$ as in
\eqref{eq:lambdaE}. Dividing by $\cV_{d_2}$ gives the formula. The rows of $\rho$
indexed by $\Gamma$ are unit rows $e_\gamma$; if $\gamma\notin E$ for some
$\gamma\in\Gamma$, that row of $\rho_{\bullet,E}$ vanishes. Otherwise expanding
along those rows leaves the minor of the $f$-band on the remaining columns.
\end{proof}

\begin{example}\label{ex:cb}
For $d_1=d_2=4$, $t=5$, $\cS=\{1,x^2\}$:
\[
\begin{aligned}
P_\cS=\;&(a_0a_3-a_1a_2)\,s_{(0,0,0,0)}+(a_0a_4-a_1a_3)\,s_{(1,0,0,0)}
        -a_1a_4\,s_{(2,0,0,0)}\\
       &+(a_3^2-a_2a_4)\,s_{(1,1,1,0)}+a_3a_4\,s_{(2,1,1,0)}
        +a_4^2\,s_{(2,2,1,0)} .
\end{aligned}
\]
The coefficients are $2\times2$ minors of the band matrix of $f$; by the dual
Jacobi--Trudi (N\"agelsbach--Kostka) identity they are, up to a power of
$a_{d_1}$ and a sign, Schur functions of the roots of $f$. Theorem~\ref{thm:skew}
makes this precise in the appropriate basis.
\end{example}

\section{The multi-Schur form and the exact sign}\label{sec:multischur}

\begin{theorem}\label{thm:multischur}
Let $\tau$ be any permutation of $\{1,\dots,d_2\}$ and let $M^\tau$ be the
$d_2\times d_2$ matrix
\[
  M^\tau_{i,q}:=X_{\varepsilon_{\tau(i)}-d_2+q}\big(A_{\tau(i)}\big),
  \qquad 1\le i,q\le d_2,
\]
where $X_r(A_i)=h_r$ if the row $\tau(i)$ is a monomial row and $X_r(A_i)=H_r$ if
it is an $f$-row, with the convention $h_r=H_r=0$ for $r<0$. Then
\[
  P_\cS=(-1)^{\binom{d_2}{2}}\,\sgn(\tau)\,\det M^\tau .
\]
In particular, taking $\tau$ to order the rows by decreasing $\varepsilon$, the
matrix $M^\tau$ is a Jacobi--Trudi matrix with one alphabet per row: the alphabet
$\bX$ on the monomial rows and the difference alphabet $\bX-\bA$ on the $f$-rows.
\end{theorem}

\begin{proof}
The key observation is that $\cO_\cS$ and $M^{\mathrm{id}}$ have \emph{the same}
coefficient matrix $\rho$, applied to two different families of rows. Indeed, put
$Y_{c,q}:=h_{c-d_2+q}(\bX)$ for $0\le c\le t^*$ and $1\le q\le d_2$. For a
monomial row of degree $\gamma$ we have $\rho_{i,\bullet}=e_\gamma$ and
$(\rho\cdot Y)_{i,q}=h_{\gamma-d_2+q}$, as required. For the $f$-row $x^if$ we
have $\rho_{i,c}=a_{c-i}$, while by \eqref{eq:hH},
\[
  H_{(i+d_1)-d_2+q}=\sum_{l=0}^{d_1}a_l\,h_{\,i+l-d_2+q}
  =\sum_{c}\rho_{i,c}\,h_{\,c-d_2+q}=(\rho\cdot Y)_{i,q},
\]
where we substituted $c=i+l$. Hence
\[
  \cO_\cS=\rho\cdot W,\qquad M^{\mathrm{id}}=\rho\cdot Y .
\]
Applying Cauchy--Binet to both, with $d_2$-subsets $E$ taken in increasing order,
\[
  |\cO_\cS|=\sum_E\det(\rho_{\bullet,E})\det(W_{E,\bullet}),\qquad
  \det M^{\mathrm{id}}=\sum_E\det(\rho_{\bullet,E})\det(Y_{E,\bullet}) .
\]
As recalled in the proof of Theorem~\ref{thm:cb},
$\det(W_{E,\bullet})=\cV_{d_2}s_{\lambda(E)}(\bX)$. On the other hand
$\det(Y_{E,\bullet})$ is a Jacobi--Trudi determinant for the same partition
$\lambda(E)$, but with the rows in the opposite order: in the Jacobi--Trudi
matrix of $\lambda(E)$ the row $p$ carries the index $e_{d_2+1-p}$, that is the
$e$'s in decreasing order, whereas $Y_{E,\bullet}$ lists them increasingly.
Reversing $d_2$ rows contributes $(-1)^{\binom{d_2}{2}}$, so
$\det(Y_{E,\bullet})=(-1)^{\binom{d_2}{2}}s_{\lambda(E)}(\bX)$. Comparing the two
sums term by term yields
$\det M^{\mathrm{id}}=(-1)^{\binom{d_2}{2}}|\cO_\cS|/\cV_{d_2}$. Finally
$\det M^\tau=\sgn(\tau)\det M^{\mathrm{id}}$, and the sign is its own inverse.
\end{proof}

\begin{remark}
The proof does not use the ordering of the rows, which is why the statement holds
for an arbitrary $\tau$: interchanging two rows changes $\sgn(\tau)$ and
$\det M^\tau$ simultaneously. In particular ties among the $\varepsilon_i$ --- which
occur whenever a monomial exponent equals $i+d_1$ --- are harmless.
\end{remark}

\section{The classical case: a rectangle}\label{sec:rect}

We now specialize to the classical scalar subresultants. By
\cite[Rmk.~2.1(2)]{DKS05}, for $0\le k\le\min\{d_1,d_2\}$ and
$t=d_1+d_2-k-1$ one takes $\cS_j:=\{x^i:0\le i\le k,\ i\ne j\}$, and then
$\Delta_{\cS_j}=(-1)^{(d_1-k)(d_2-k)}S_k^{(j)}$, where $S_k^{(j)}$ is the scalar
subresultant of \cite[(3)]{DKS05}. Note that $t\ge d_2-1$, so $t^*=t$ and there
are no rows $x^{t+1},\dots,x^{t^*}$; moreover $\Gamma=\{0,\dots,k\}\smallsetminus\{j\}$
and $\Phi=\{d_1,\dots,t\}$ has $m:=d_2-k$ elements.

\begin{theorem}\label{thm:rect}
Let $\cS_k=\{1,x,\dots,x^{k-1}\}$, i.e.\ $j=k$, and $m=d_2-k$. Then
\[
  P_{\cS_k}
  =\det\big(H_{(d_1-k)-p+q}\big)_{p,q=1}^{m}
  =a_{d_1}^{\,d_2-k}\;s_{\big((d_1-k)^{d_2-k}\big)}(\bX-\bA),
\]
with sign $+1$: the supersymmetric Schur function of the rectangular partition
$(d_1-k)\times(d_2-k)$ in the difference alphabet.
\end{theorem}

\begin{proof}
Here $\Gamma=\{0,1,\dots,k-1\}$ and the degrees are
$\varepsilon=(0,1,\dots,k-1,d_1,d_1+1,\dots,d_1+m-1)$, already increasing.
Ordering by decreasing $\varepsilon$ is therefore the full reversal of $d_2$ rows,
whose sign $(-1)^{\binom{d_2}{2}}$ cancels the one in
Theorem~\ref{thm:multischur}; so $P_{\cS_k}=\det M^\tau$ with $M^\tau$ the
Jacobi--Trudi matrix in decreasing order. Write $M^\tau$ in blocks, with the $m$
$f$-rows first (row $p=1,\dots,m$ carrying $\varepsilon=d_1+m-p$) and the $k$
monomial rows below (row $p'=1,\dots,k$ carrying $\gamma=k-p'$), and split the
columns as $q\le m$ and $q=m+r$ with $1\le r\le k$.

\emph{Lower left block ($k\times m$) vanishes.} Its entry is
$h_{\gamma-d_2+q}$ with $\gamma\le k-1$ and $q\le m=d_2-k$, so
$\gamma-d_2+q\le-1<0$.

\emph{Lower right block ($k\times k$) is unitriangular.} With $q=m+r$ its entry is
$h_{(k-p')-d_2+(d_2-k+r)}=h_{r-p'}$, which is upper triangular with $h_0=1$ on the
diagonal; its determinant is $1$.

\emph{Upper left block ($m\times m$).} Its entry is
$H_{(d_1+m-p)-d_2+q}=H_{(d_1-k)-p+q}$, because $m=d_2-k$.

Being block upper triangular, $\det M^\tau$ is the product of the two diagonal
blocks, i.e.\ the Jacobi--Trudi determinant of the rectangle $((d_1-k)^m)$ in the
$H$'s. By \eqref{eq:hH} each of the $m$ rows carries a factor $a_{d_1}$, giving
the stated normalization.
\end{proof}

\begin{corollary}\label{cor:LP}
$S_k^{(k)}=(-1)^{(d_1-k)(d_2-k)}\,b_{d_2}^{\,d_1-k}a_{d_1}^{\,d_2-k}\,
 s_{((d_1-k)^{d_2-k})}(\bX-\bA)$.
\end{corollary}

\begin{proof}
Combine Theorem~\ref{thm:rect} with \eqref{eq:dks} and
$\Delta_{\cS_k}=(-1)^{(d_1-k)(d_2-k)}S_k^{(k)}$, using $t^*-d_2+1=d_1-k$ and
$\sg(\cS_k)=+1$; the latter because for $\cS_k=\{1,\dots,x^{k-1}\}$ with $t^*=t$
the reordering defining $\sg$ in \cite[\S2]{DKS05} is the identity.
\end{proof}

Corollary~\ref{cor:LP} is the statement of Lascoux and Pragacz \cite{LP} for the
principal subresultant coefficient, here obtained from \eqref{eq:dks}. The
extreme cases are consistent: $k=0$ gives the rectangle $d_1\times d_2$, i.e.\ the
resultant, and $k=\min\{d_1,d_2\}$ gives the empty partition.

\begin{remark}\label{rmk:gaps}
The proof uses in an essential way that the monomial exponents
$\Gamma=\{0,\dots,k-1\}$ are \emph{consecutive}: this is what makes the lower left
block vanish and the lower right block unitriangular. For $j<k$ the set
$\{0,\dots,k\}\smallsetminus\{j\}$ has a gap and the collapse fails. For instance
for $d_1=d_2=3$, $k=2$, $j=0$ one gets
\[
  P_{\cS_0}=a_0+a_3\,e_3(\bX)=a_3\big(e_3(\bX)-e_3(\bA)\big),
\]
which is the mixed determinant
\[
  \det\begin{pmatrix}H_1&H_2&H_3\\ 1&h_1&h_2\\ 0&1&h_1\end{pmatrix},
\]
that is, the truncated alternating sum $\sum_{i=0}^{2}(-1)^iH_{3-i}e_i(\bX)$;
the untruncated sum would collapse to $[z^3]f^\rev(z)=a_0$. An exhaustive search
over partitions confirms that $P_{\cS_j}$ is not a single supersymmetric Schur
function when $j<k$.
\end{remark}

\section{Toeplitz determinants and Pad\'e approximation}\label{sec:toeplitz}

\begin{theorem}\label{thm:toeplitz}
With $c_r$ as in \eqref{eq:c} and $m=d_2-k$,
\[
  S_k^{(k)}=(-1)^{(d_1-k)(d_2-k)}\,b_{d_2}^{\,d_1+d_2-2k}\,
  \det\big(c_{(d_1-k)-p+q}\big)_{p,q=1}^{m}.
\]
\end{theorem}

\begin{proof}
By \eqref{eq:c}, $H_r=b_{d_2}c_r$, so the $m\times m$ determinant of
Theorem~\ref{thm:rect} equals $b_{d_2}^{m}\det(c_{(d_1-k)-p+q})$. Now apply
Corollary~\ref{cor:LP}, or equivalently multiply by $b_{d_2}^{d_1-k}$ as in its
proof; $m+(d_1-k)=d_1+d_2-2k$.
\end{proof}

Since $c_r$ are the Taylor coefficients of $f^\rev/g^\rev$, i.e.\ of the expansion
of $f/g$ at infinity, Theorem~\ref{thm:toeplitz} says that the principal
subresultant coefficient is a Toeplitz determinant of that series. This is exactly
the determinantal criterion underlying the Pad\'e approximation problem, and it
makes explicit the link with the rational interpolation viewpoint of
\cite{DKS15}.

\section{The natural basis for arbitrary $\cS$}\label{sec:skew}

By Remark~\ref{rmk:gaps} the quotient is not a single supersymmetric Schur
function in general. It does, however, expand in the basis
$\{s_\nu(\bX-\bA)\}$, and the expansion is dual to that of Theorem~\ref{thm:cb}.
Throughout this section we normalize $f$ to be monic, $a_{d_1}=1$, so that the
$h_s(\bA)$ are polynomials in the $a_l$; in general they lie in
$\ZZ[a][a_{d_1}^{-1}]$.

Set $u_r:=h_r(\bX-\bA)$ and $U_{r,q}:=u_{r-d_2+q}$.

\begin{theorem}\label{thm:skew}
Let $\cC$ be the matrix whose rows are indexed as in Definition~\ref{def:rows},
with
\[
  \cC_{i,c}=\delta_{c,\varepsilon_i}\ \ \text{for an $f$-row},\qquad
  \cC_{i,c}=h_{\gamma-c}(\bA)\ \ \text{for a monomial row of degree }\gamma .
\]
Then $M^{\mathrm{id}}=\cC\cdot U$ and consequently
\[
  P_\cS=(-1)^{\binom{d_2}{2}}\sum_{E}\det\big(\cC_{\bullet,E}\big)\,
        \det\big(U_{E,\bullet}\big)
      =\sum_{F}\pm\,\det\big(\cC_{\bullet,\Phi\sqcup F}\big)\,
        s_{\nu(\Phi\sqcup F)}(\bX-\bA),
\]
where $F$ runs over the $\#\Gamma$-subsets of $\ZZ_{\ge0}\smallsetminus\Phi$ and
$\nu(\cdot)$ is as in \eqref{eq:lambdaE}. Moreover, taking the rows (indexed by
$\Gamma$) and the columns (indexed by $F$) in decreasing order,
\[
  \det\big(\cC_{\bullet,\Phi\sqcup F}\big)
  =\pm\det\big(h_{\gamma_p-f_q}(\bA)\big)_{p,q}
  =\pm\,s_{\lambda(\Gamma)/\mu(F)}(\bA),
  \qquad
  \lambda(\Gamma)_p=\gamma_p+p,\quad \mu(F)_q=f_q+q .
\]
\end{theorem}

\begin{proof}
Since $\bX=(\bX-\bA)+\bA$ we have $h_r(\bX)=\sum_s h_s(\bA)h_{r-s}(\bX-\bA)$,
which is exactly the monomial row of $\cC$ applied to $U$; and
$H_r=a_{d_1}u_r=u_r$ for monic $f$, which is the unit row. Hence
$M^{\mathrm{id}}=\cC\cdot U$, and Cauchy--Binet together with
Theorem~\ref{thm:multischur} gives the first display, the two occurrences of
$(-1)^{\binom{d_2}{2}}$ cancelling exactly as in the proof of
Theorem~\ref{thm:multischur}. The unit rows force $\Phi\subseteq E$, whence the
reindexing $E=\Phi\sqcup F$. Finally, the surviving minor is
$\det(h_{\gamma_p-f_q}(\bA))$, and with both indices decreasing the sequences
$\lambda_p=\gamma_p+p$ and $\mu_q=f_q+q$ are weakly decreasing, so this is the
skew Jacobi--Trudi determinant of $\lambda(\Gamma)/\mu(F)$, cf.\
\cite[I.(5.4)]{Macdonald}.
\end{proof}

\begin{corollary}\label{cor:support}
$\det(\cC_{\bullet,\Phi\sqcup F})\neq0$ if and only if
$\mu(F)\subseteq\lambda(\Gamma)$ and the skew diagram
$\lambda(\Gamma)/\mu(F)$ has no column of length greater than $d_1$.
\end{corollary}

\begin{proof}
$s_{\lambda/\mu}=0$ unless $\mu\subseteq\lambda$, and a skew Schur function
vanishes on an alphabet of $d_1$ letters exactly when the skew diagram has a
column of length $>d_1$; see \cite[I.5]{Macdonald}. Here $\#\bA=d_1$.
\end{proof}

Corollary~\ref{cor:support} explains the sparsity of the expansion and makes
precise the role of the gaps of $\Gamma$: for $\Gamma=\{k-1,\dots,0\}$ one gets
$\lambda(\Gamma)_p=(k-p)+p=k$, i.e.\ $\lambda$ is the $k\times k$ square, and a
single $F$ survives, recovering Theorem~\ref{thm:rect}.

\begin{example}
With $f$ monic, the expansions are short. All Schur functions below are taken in
the alphabet $\bX-\bA$:
\[
\begin{array}{r@{\quad}l}
  d_1=3,\ d_2=3,\ t=4,\ \cS=\{x\}:      & P_\cS=-a_2\,s_{(2,2)}+s_{(2,2,1)},\\[2pt]
  d_1=2,\ d_2=3,\ t=3,\ \cS=\{x^2\}:    & P_\cS=(a_1^2-a_0)\,s_{(1,1)}-a_1\,s_{(1,1,1)},\\[2pt]
  d_1=4,\ d_2=4,\ t=5,\ \cS=\{1,x^2\}:  & P_\cS=-a_3\,s_{(2,2)}+s_{(2,2,1)},\\[2pt]
  d_1=3,\ d_2=3,\ t=3,\ \cS=\{1,x\}:    & P_\cS=-s_{(1)} .
\end{array}
\]
The last line has a single term, as it must: there $\Gamma=\{0,1\}$ is an
interval. The coefficients are the $h_s(\bA)$: for monic $f$,
$h_1(\bA)=-a_{d_1-1}$, $h_2(\bA)=a_{d_1-1}^2-a_{d_1-2}$, and so on.
\end{example}

\section{Multiple roots: the Newton form}\label{sec:newton}

If $g$ has repeated roots both $|\cO_\cS|$ and $\cV_{d_2}$ vanish and
\eqref{eq:dks} is a $0/0$ expression. The following identity removes the
difficulty at once. Recall that the divided difference of a polynomial $R$ at
nodes $\xi_1,\dots,\xi_j$ is
$R[\xi_1,\dots,\xi_j]=\sum_{i=1}^{j}R(\xi_i)\big/\prod_{l\ne i}(\xi_i-\xi_l)$.

\begin{theorem}\label{thm:newton}
With the row polynomials $R_i$ of Definition~\ref{def:rows},
\[
  P_\cS=\det\big(R_i[\xi_1,\dots,\xi_j]\big)_{1\le i,j\le d_2}.
\]
\end{theorem}

\begin{proof}
Let $T$ be the $d_2\times d_2$ matrix expressing the divided differences in terms
of the evaluations, $T_{j,j'}=1/\prod_{l\le j,\,l\ne j'}(\xi_{j'}-\xi_l)$ for
$j'\le j$ and $T_{j,j'}=0$ for $j'>j$. It is lower triangular with diagonal
entries $T_{j,j}=1/\prod_{l<j}(\xi_j-\xi_l)$, so
$\det T=\prod_{l<j}(\xi_j-\xi_l)^{-1}=\cV_{d_2}^{-1}$. Since
$\big(R_i[\xi_1,\dots,\xi_j]\big)_{i,j}=\cO_\cS\cdot T^{\mathsf t}$, taking
determinants gives the claim.
\end{proof}

\begin{corollary}\label{cor:conf}
Let $g$ have distinct roots $y_1,\dots,y_s$ with multiplicities $m_1,\dots,m_s$,
$\sum_i m_i=d_2$. Let $\cO^{\mathrm{conf}}_\cS$ be the confluent matrix whose
columns are indexed by the pairs $(i,r)$, $0\le r<m_i$, with entries
$\frac{1}{r!}R^{(r)}(y_i)$, and let
$\cV^{\mathrm{conf}}=\prod_{i<i'}(y_{i'}-y_i)^{m_im_{i'}}$. Then
\[
  \frac{|\cO^{\mathrm{conf}}_\cS|}{\cV^{\mathrm{conf}}}
  =P_\cS\big(y_1,\dots,y_1,\ \dots,\ y_s,\dots,y_s\big),
\]
each $y_i$ repeated $m_i$ times. In particular the case of multiple roots needs no
new formula: the universal polynomial $P_\cS$ evaluated at the multiset of roots
computes it.
\end{corollary}

\begin{proof}
Split the roots into $s$ groups, the $i$-th consisting of
$\xi_{i,1},\dots,\xi_{i,m_i}$, and let $\cV^{(i)}:=\prod_{l<j}(\xi_{i,j}-\xi_{i,l})$
be the Vandermonde determinant of the $i$-th group. Performing inside each group
the column operations of Theorem~\ref{thm:newton}, which replace the columns
$R(\xi_{i,1}),\dots,R(\xi_{i,m_i})$ by the divided differences
\[
  R[\xi_{i,1}],\quad \dots,\quad R[\xi_{i,1},\dots,\xi_{i,m_i}],
\]
multiplies the determinant by $\prod_i(\cV^{(i)})^{-1}$. Hence
\[
  \frac{|\cO_\cS|}{\prod_i \cV^{(i)}}
  =\det\Big(\text{within-group divided differences}\Big),
\]
and on the other hand, splitting the factors of $\cV_{d_2}$ according to whether
the two roots lie in the same group or not,
\[
  \frac{\cV_{d_2}}{\prod_i \cV^{(i)}}
  =\prod_{i<i'}\ \prod_{j,j'}\big(\xi_{i',j'}-\xi_{i,j}\big).
\]
Now let $\xi_{i,j}\to y_i$ for every $j$. The left-hand sides converge to
$|\cO^{\mathrm{conf}}_\cS|$ and $\cV^{\mathrm{conf}}$ respectively, because
divided differences at coalescing nodes converge to the divided derivatives,
$R[\underbrace{y,\dots,y}_{r+1}]=R^{(r)}(y)/r!$, and because the second product
converges to $\prod_{i<i'}(y_{i'}-y_i)^{m_im_{i'}}$. Since
$P_\cS=|\cO_\cS|/\cV_{d_2}$ is the ratio of the two left-hand sides and is a
polynomial, its value at the limit is the limit of the values, which gives the
statement.
\end{proof}

\section{The multivariate case: an obstruction}\label{sec:multi}

In the multivariate setting \cite[Thm.~3.2]{DKS05} the denominator is the
generalized Vandermonde determinant $\cV_T=\det(\xi_j^{\alpha_i})$ built on the
$\mathbf d=d_1\cdots d_n$ common roots $\xi_j\in\overline K^{\,n}$ of
$f_1,\dots,f_n$. The analogue of Theorem~\ref{thm:div} fails.

\begin{proposition}\label{prop:multi}
Let $n=2$, $d_1=d_2=d_3=2$ and $t=2$ as in \cite[Ex.~3.4]{DKS05}, so that
\[
  |\cO_\cS|=-c_0|V_{2,5}|+c_2|V_{0,5}|+c_5|V_{0,2}|,\qquad \cV_T=|V_{4,5}|,
\]
where $V$ is the $6\times4$ matrix of evaluations of $1,u,v,uv,u^2,v^2$ at the
four roots $\xi_j=(u_j,v_j)$ and $V_{i,i'}$ denotes the $4\times4$ submatrix
obtained by deleting rows $i,i'$. Then $\cV_T$ divides none of $|V_{2,5}|$,
$|V_{0,5}|$, $|V_{0,2}|$ in $\QQ[u_1,\dots,u_4,v_1,\dots,v_4]$, and hence
$\cV_T\nmid|\cO_\cS|$.
\end{proposition}

\begin{proof}[Proof (computational)]
The coefficients $c_0,c_2,c_5$ of $f_3$ are independent indeterminates, so
$\cV_T$ divides $|\cO_\cS|$ if and only if it divides each of the three
alternants separately. A direct division in $\QQ[u,v]$ shows that it divides none
of them; in fact none of the $14$ determinants $|V_{i,i'}|$ with
$(i,i')\ne(4,5)$ is divisible by $\cV_T$. The test is conclusive because here the
roots are free parameters: four general points of the plane are the base locus of
a pencil of conics, so the map from coefficients to roots is dominant and the
eight coordinates $u_j,v_j$ are algebraically independent; therefore
``divisible in $\QQ[u,v]$'' is the same as ``universally polynomial in the
roots''. The degrees are $\deg\cV_T=4$ and $\deg|V_{2,5}|=5$,
$\deg|V_{0,5}|=6$, $\deg|V_{0,2}|=7$.
\end{proof}

The reason is structural. The proof of Theorem~\ref{thm:div} rests on
Lemma~\ref{lem:alt}, i.e.\ on the fact that in $K[\xi_1,\dots,\xi_d]$ with scalar
$\xi_j$ the module of alternating polynomials is free of rank one over the
symmetric ones, generated by the Vandermonde determinant. For points of
$\mathbb{A}^n$ with $n\ge2$ the corresponding module of diagonally alternating
polynomials is not principal --- this is the setting of diagonal harmonics, cf.\
\cite{Haiman} --- and $\cV_T$ is merely one of several generators. Consequently
the univariate statement cannot be transcribed, and the multivariate problem must
be reformulated. Natural reformulations are: to exhibit explicitly the
cancellation of $\cV_T$ against the product $\prod_j\overline\Delta_{T_j}$ of
homogeneous subresultants appearing in \cite[(13)]{DKS05}, which \emph{is}
polynomial in the coefficients; to determine the minimal universal denominator,
i.e.\ $\gcd(|\cO_\cS|,\cV_T)$; or to expand $|\cO_\cS|$ in a generating set of the
module of alternating polynomials. We leave this for future work.

\section{Symbolic verification}\label{sec:verif}

All statements above were verified by symbolic computation, with two independent
implementations (SymPy 1.14 and SageMath 10.7); the Schur expansions were
cross-checked against the native symmetric function routines of SageMath. We
record the scope.

\begin{center}
\renewcommand{\arraystretch}{1.25}
\begin{tabular}{>{\raggedright\arraybackslash}p{0.30\textwidth}%
                >{\raggedright\arraybackslash}p{0.60\textwidth}}
\hline
Statement & Range verified\\
\hline
Thms.~\ref{thm:div}, \ref{thm:cb}, \ref{thm:multischur}
  & $12$ configurations $(d_1,d_2,t,\cS)$ with $d_1,d_2\le6$, including
    $t<d_2-1$ and ties among the $\varepsilon_i$; among them
    \cite[Ex.~2.3 and 2.4]{DKS05}.\\
Thm.~\ref{thm:rect}
  & all $1\le d_1,d_2\le5$ and all $0\le k\le\min\{d_1,d_2\}$, without search.\\
Thm.~\ref{thm:toeplitz}
  & $46$ cases, $1\le d_1,d_2\le4$ and all $k$, with $a$ and $b$ independent
    (not passing through the roots).\\
Thm.~\ref{thm:skew}, Cor.~\ref{cor:support}
  & $5$ configurations; the vanishing criterion was tested on \emph{all} $36$
    admissible subsets $F$.\\
Thm.~\ref{thm:newton}
  & $6$ configurations; Cor.~\ref{cor:conf} on $8$ multiplicity patterns
    $[2],[2,1],[3],[2,2],[3,1],[2,1,1]$.\\
Prop.~\ref{prop:multi}
  & all $15$ alternants of \cite[Ex.~3.4]{DKS05}.\\
\hline
\end{tabular}
\end{center}

Finally, everything above is relative to \eqref{eq:dks}. As an end-to-end check we
also verified, for $20$ tuples $(d_1,d_2,k,j)$, the identity
\[
  S_k^{(j)}=(-1)^{(d_1-k)(d_2-k)}\,\sg(\cS_j)\,b_{d_2}^{\,d_1-k}\,P_{\cS_j}
\]
with $S_k^{(j)}$ computed from its classical determinantal definition
\cite[(3)]{DKS05} and $b_l$ rewritten in the roots as
$b_l=b_{d_2}(-1)^{d_2-l}e_{d_2-l}(\bX)$, \emph{with the exact signs}, including
$\sg(\cS_j)$. This is the only check that does not take \eqref{eq:dks} as a
hypothesis, and it validates it as well.

% ---------------------------------------------------------------------------
% TODO: check the exact bibliographic data of [Hong], [DKS13], [DHKS],
% [DKS15] before circulating the manuscript.
% ---------------------------------------------------------------------------
\begin{thebibliography}{99}

\bibitem{DKS05} C.~D'Andrea, T.~Krick, A.~Szanto,
  \emph{Multivariate subresultants in roots},
  J.~Algebra \textbf{302} (2006), 16--36.

\bibitem{DKS13} C.~D'Andrea, T.~Krick, A.~Szanto,
  \emph{Subresultants in multiple roots},
  Linear Algebra Appl.\ \textbf{438} (2013).

\bibitem{DKS15} C.~D'Andrea, T.~Krick, A.~Szanto,
  \emph{Subresultants, Sylvester sums and the rational interpolation problem},
  J.~Symbolic Comput.\ \textbf{68} (2015).

\bibitem{DHKS} C.~D'Andrea, H.~Hong, T.~Krick, A.~Szanto,
  \emph{An elementary proof of Sylvester's double sums for subresultants},
  J.~Symbolic Comput.\ \textbf{42} (2007), 290--297; and
  \emph{Sylvester's double sums: the general case},
  J.~Symbolic Comput.\ \textbf{44} (2009).

\bibitem{Hong} H.~Hong, \emph{Subresultants in roots}, technical report,
  North Carolina State University.

\bibitem{LP} A.~Lascoux, P.~Pragacz,
  \emph{Double Sylvester sums for subresultants and multi-Schur functions},
  J.~Symbolic Comput.\ \textbf{35} (2003), 689--710.

\bibitem{Lascoux} A.~Lascoux,
  \emph{Symmetric Functions and Combinatorial Operators on Polynomials},
  CBMS Regional Conference Series in Mathematics \textbf{99}, AMS, 2003.

\bibitem{Macdonald} I.~G.~Macdonald,
  \emph{Symmetric Functions and Hall Polynomials}, 2nd ed.,
  Oxford University Press, 1995.

\bibitem{Chardin} M.~Chardin, \emph{Multivariate subresultants},
  J.~Pure Appl.\ Algebra \textbf{101} (1995), 129--138.

\bibitem{Macaulay} F.~S.~Macaulay, \emph{Some formulae in elimination},
  Proc.\ London Math.\ Soc.\ \textbf{35} (1902), 3--27.

\bibitem{Haiman} M.~Haiman,
  \emph{Conjectures on the quotient ring by diagonal invariants},
  J.~Algebraic Combin.\ \textbf{3} (1994), 17--76.

\end{thebibliography}

\end{document}
